\chapter{Introduction}
\label{cha:Introduction}

While hardware and software are important for modern audio signal processing, software has established itself as the more economical tool due to its flexibility, scalability and lower costs. Thus, we find increasingly innovative applications through findings in digital signal processing (DSP) and the advancement of modern technology, such as for sound enhancement in real-time applications for video conferencing, audio signal processing algorithms are used for noise cancellation, echo cancellation. Digital audio workstations are used in music production to process and mix recorded audio signals. They also offer a variety of audio plug-ins that virtually recreate real instruments. In view of the latest developments in deep learning, a big step has been taken in automatic speech recognition, speech separation, audio optimisation and synthesising ambient noise.  

This involves converting signals, such as audio or video, from analogue to digital so that the signal data can be analysed and manipulated by the computer using mathematical algorithms.  

\section{Objectives}

The aim of the thesis is to carry out a comprehensive evaluation of the FAUST programming language by illustrating its advantages and disadvantages in comparison to other programming languages in the field of audio signal processing. Properties such as performance and flexibility of FAUST are analysed and evaluated in detail. In addition, the technological possibilities offered by FAUST will be discussed.  


\section{Why FAUST?}

Technological progress has led to the development of a number of specialised programming languages and architectures to enable more efficient and flexible processing of audio signals. A wide range of programming languages are now available, ranging from machine-friendly low-level languages to programmer-friendly high-level languages.  

FAUST (Functional Audio Stream), on the other hand, has the advantages of both worlds at the same time. With a very intuitive programming environment and its own compiler, FAUST enables its code to be translated into an equivalent programme in languages such as C++, Java, LLVM, WebAssembly and more. Compared to the C++ code of experienced developers, the FAUST generated code can both keep up and even outperform the developers.  

FAUST is characterised by its good documentation and beginner-friendliness, so that no advanced DSP knowledge is required. Nevertheless, the basic concepts of DSP make it easier to understand and use FAUST. The following chapters cover the basics of DSP and explain the features of FAUST. In addition, the functionalities, performance and technological possibilities of FAUST are explored.  

